\subsection{PERMUTABLE ELEMENTS. COMMUTATIVE LAWS}

\begin{enumerate}
\def\labelenumi{(\arabic{enumi})}
\setcounter{enumi}{3}
\tightlist
\item
  Let \((x, y) \mapsto x \top y\) be a commutative law on E; the law
\end{enumerate}

\[
(\mathbf {X}, \mathbf {Y}) \mapsto \mathbf {X} \top \mathbf {Y}
\]

between subsets of E is commutative.

DEFINITION 9. Let E be a magma and X a subset of E. The set of elements of E which commute with each of the elements of X is called the centralizer of X.

Let X and Y be two subsets of E and X' and Y' their respective centralizers. If X ⊂ Y, then Y' ⊂ X'.

Let \((\mathbf{X}_i)_{i\in I}\) be a family of subsets of \(E\) and for all \(i\in I\) let \(\mathbf{X}_i'\) be the centralizer of \(\mathbf{X}_i\) . The centralizer of \(\bigcup_{i\in I}\mathbf{X}_i\) is \(\bigcap_{i\in I}\mathbf{X}_i'\) .

Let X be a subset of E and \(X'\) the centralizer of X. The centralizer \(X''\) of \(X'\) is called the bicentralizer of X. Then \(X \subset X''\) . The centralizer \(X'''\) of \(X''\) is equal to \(X'\) . For \(X'\) is contained in its bicentralizer \(X'''\) and the relation \(X \subset X''\) implies \(X''' \subset X'\) .

PROPOSITION 3. Let E be an associative magma whose law is denoted by \(\top\) . If an element \(x\) of E commutes with each of the elements \(y\) and \(z\) of E, it commutes with \(y \top z\) .

For

\[
\begin{array}{r l} x \top (y \top z) = (x \top y) \top z & = (y \top x) \top z \\ & = y \top (x \top z) = y \top (z \top x) = (y \top z) \top x. \end{array}
\]

COROLLARY. Let E be an associative magma. The centralizer of any subset of E is a stable subset of E.

DEFINITION 10. The centralizer of a magma E is called the centre of E. An element of the centre of E is called a central element of E.

If E is an associative magma its centre is a stable subset by the Corollary to Proposition 3 and the law induced on its centre is commutative.

PROPOSITION 4. Let E be an associative magma, X and Y two subsets of E. If every element of X commutes with every element of Y every element of the stable subset generated by X commutes with every element of the stable subset generated by Y.

Let \(X'\) and \(X''\) be the centralizer and bicentralizer of X. They are stable subsets of E. Now \(X \subset X''\) and \(Y \subset X'\) and hence \(X''\) (resp. \(X'\) ) contains the stable subset of E generated by X (resp. Y). As every element of \(X''\) commutes with every element of \(X'\) , the proposition follows.

COROLLARY 1. If \(x\) and \(y\) are permutable under the associative law \(\top\) so are \(\top^m x\) and \(\top^n y\) for all integers \(m > 0\) and \(n > 0\) ; in particular \(\top^m x\) and \(\top^n x\) are permutable for all \(x\) and all integers \(m > 0\) and \(n > 0\) .

COROLLARY 2. If all pairs of elements of a subset X are permutable under an associative law T, the law induced by T on the stable subset generated by X is associative and commutative.

THEOREM 2 (Commutativity theorem). Let \(\top\) be an associative law of composition on E; let \((x_{\alpha})_{\alpha \in A}\) be a non-empty finite family of elements of E which are pairwise permutable; let B and C be two totally ordered sets with A as underlying set. Then \(\bigcap_{\alpha \in B} x_{\alpha} = \bigcap_{\alpha \in C} x_{\alpha}\) .

Since the theorem is true if A has a single element, we argue by induction on the number p of elements in A. Let p be an integer \textgreater1 and suppose the theorem is true when Card A \textless{} p.~We prove it for Card A = p.~It may be assumed that A is the interval \([0, p - 1]\) in N; the composition of the ordered sequence \((x_{\alpha})_{\alpha \in \mathbf{A}}\) defined by the natural order relation on \(\mathbf{A}\) is \(\prod_{i = 0}^{p - 1}x_i\) .

Let A be given another total ordering and let h be the least element of A under this ordering and \(A'\) the set of other elements of A (totally ordered by the induced ordering). Suppose first 0 \textless{} h \textless{} p - 1 and let \(P = \{0, 1, \ldots, h - 1\}\) and \(Q = \{h + 1, \ldots, p - 1\}\) ; the theorem being assumed true for \(A'\) , applying the associativity theorem, we obtain (since \(A' = P \cup Q\) )

\[
\bigcap_ {\alpha \in \mathbf {A} ^ {\prime}} x _ {\alpha} = \left(\prod_ {i = 0} ^ {h - 1} x _ {i}\right) \top \left(\prod_ {i = h + 1} ^ {p - 1} x _ {i}\right)
\]

whence, composing \(x_{h}\) with both sides and repeatedly applying the commutativity and associativity of T:

\[
\begin{array}{r l} \underset {\alpha \in \mathbf {A}} {\top} x _ {\alpha} & = x _ {h} \top \left(\underset {\alpha \in \mathbf {A} ^ {\prime}} {\top} x _ {\alpha}\right) = x _ {h} \top \left(\underset {i = 0} {\overset {h - 1} {\top}} x _ {i}\right) \top \left(\underset {i = h + 1} {\overset {p - 1} {\top}} x _ {i}\right) \\ & = \left(\underset {i = 0} {\overset {h - 1} {\top}} x _ {i}\right) \top x _ {h} \top \left(\underset {i = h + 1} {\overset {p - 1} {\top}} x _ {i}\right) = \underset {i = 0} {\overset {p - 1} {\top}} x _ {i}, \end{array}
\]

which proves the theorem for this case. If h = 0 or h = p - 1, the same result follows, but more simply, the terms arising from P or the terms arising from Q not appearing in the formulae.

Under a commutative associative law on a set E the composition of a finite family \((x_{\alpha})_{\alpha \in A}\) of elements of E is by definition the common value of the composition of all the ordered sequences obtained by totally ordering A in all possible ways. This composition will still be denoted by \(\bigcap_{\alpha \in A} x_{\alpha}\) under a law denoted by T; similarly for other notations.

THEOREM 3. Let \(\top\) be an associative law on E and \((x_{\alpha})_{\alpha \in \mathbf{A}}\) a non-empty finite family of elements of E which are pairwise permutable. If A is a union of non-empty subsets \((\mathbf{B}_i)_{i\in I}\) which are pairwise disjoint, then

\[
\bigcap_ {\alpha \in A} x _ {\alpha} = \bigcap_ {i \in I} \left(\bigcap_ {\alpha \in B _ {i}} x _ {\alpha}\right).\tag{6}
\]

This follows from Theorem 2 if A and I are totally ordered so that the \(B_{i}\) satisfy the conditions of Theorem 1.

We single out two important special cases of this theorem:

\begin{enumerate}
\def\labelenumi{\arabic{enumi}.}
\tightlist
\item
  If \((x_{\alpha\beta})_{(\alpha,\beta)\in\mathbf{A}\times\mathbf{B}}\) is a finite family of permutable elements of an associative magma whose indexing set is the product of two non-empty finite sets A, B (a ``double family''), then
\end{enumerate}

\[
\underset {(\alpha , \beta) \in A \times B} {\top} x _ {\alpha \beta} = \underset {\alpha \in A} {\top} \left(\underset {\beta \in B} {\top} x _ {\alpha \beta}\right) = \underset {\beta \in B} {\top} \left(\underset {\alpha \in A} {\top} x _ {\alpha \beta}\right)\tag{7}
\]

as follows from Theorem 3 by considering \(\mathbf{A} \times \mathbf{B}\) as the union of the sets \(\{\alpha\} \times \mathbf{B}\) on the one hand and of the sets \(\mathbf{A} \times \{\beta\}\) on the other.

For example, if B has n elements and for each \(\alpha \in A\) all the \(x_{\alpha\beta}\) have the same value \(x_{\alpha}\) , then

\[
\bigcap_ {\alpha \in \mathbf {A}} \left(\frac {n}{T} x _ {\alpha}\right) = \frac {n}{T} \left(\bigcap_ {\alpha \in \mathbf {A}} x _ {\alpha}\right).\tag{8}
\]

If B has two elements, we obtain the following results: let \((x_{\alpha})_{\alpha \in A}, (y_{\alpha})_{\alpha \in A}\) be two non-empty families of elements of E. If the \(x_{\alpha}\) and the \(y_{\beta}\) are pairwise permutable, then

\[
\underset {\alpha \in \mathbf {A}} {\top} x _ {\alpha} \top y _ {\alpha} = \left(\underset {\alpha \in \mathbf {A}} {\top} x _ {\alpha}\right) \top \left(\underset {\alpha \in \mathbf {A}} {\top} y _ {\alpha}\right).\tag{9}
\]

Because of formula (7) the composition of a double sequence \((x_{ij})\) whose indexing set is the product of two intervals \(\{p, q\}\) and \(\{r, s\}\) in \(\mathbf{N}\) is often denoted for a commutative associative law written additively by

\[
\sum_ {i = p} ^ {q} \sum_ {j = r} ^ {s} x _ {i j} \quad \text { or } \quad \sum_ {j = r} ^ {s} \sum_ {i = p} ^ {q} x _ {i j}
\]

and similarly for laws denoted by other symbols.

\begin{enumerate}
\def\labelenumi{\arabic{enumi}.}
\setcounter{enumi}{1}
\tightlist
\item
  Let \(n\) be an integer \(>0\) and let \(A\) be the set of ordered pairs of integers \((i,j)\) such that \(0 \leqslant i \leqslant n, 0 \leqslant j \leqslant n\) and \(i < j\) ; the composition of a family \((x_{ij})_{(i,j) \in A}\) (under a commutative associative law) is also denoted by \(\bigcap_{0 \leqslant i < j \leqslant n} x_{ij}\) (or simply \(\bigcap_{i < j} x_{ij}\) , if no confusion arises); Theorem 3 here gives the formulae
\end{enumerate}

\[
\underset {0 \leqslant i <   j \leqslant n} {\mathsf {T}} x _ {i j} = \underset {i = 0} {\overset {n - 1} {\mathsf {T}}} \left(\underset {j = i + 1} {\overset {n} {\mathsf {T}}} x _ {i j}\right) = \underset {j = 1} {\overset {n} {\mathsf {T}}} \left(\underset {i = 0} {\overset {j - 1} {\mathsf {T}}} x _ {i j}\right).\tag{10}
\]

There are analogous formulae to (7) for a family whose indexing set is the product of more than two sets and analogous formulae to (10) for a family whose indexing set is the set \(S_{p}\) of strictly increasing sequences \((i_{k})_{1\leqslant k\leqslant p}\) of p integers such that \(0\leqslant i_{k}\leqslant n\) ( \(p\leqslant n+1\) ): in the latter case the composition of the family \((x_{i_{1}i_{2}\ldots i_{p}})(i_{1},\ldots,i_{p})\in\mathbb{S}_{p}\) is denoted by

\[
\underset {0 \leqslant i _ {1} <   i _ {2} <   \dots <   i _ {p} <   n} {\mathsf {T}} x _ {i _ {1} i _ {2} \ldots i _ {p}}, \quad \text {or simply} \quad \underset {i _ {1} <   i _ {2} <   \dots <   i _ {p}} {\mathsf {T}} x _ {i _ {1} i _ {2} \ldots i _ {p}}.
\]

PROPOSITION 5. Let E and F be magmas whose laws are denoted by \(\top\) and let \(f\) and \(g\) be homomorphisms of E into F. Let \(f \top g\) be the mapping \(x \mapsto f(x) \top g(x)\) of E into F. If F is associative and commutative, \(f \top g\) is a homomorphism.

For all elements \(x\) and \(y\) of \(\mathbf{E}\) :

\[
\begin{array}{r l} (f \top g) (x \top y) & = f (x \top y) \top g (x \top y) = f (x) \top f (y) \top g (x) \top g (y) \\ & = f (x) \top g (x) \top f (y) \top g (y) = ((f \top g) (x)) \top ((f \top g) (y)). \end{array}
\]

